% NCD-CIE Clinician Evaluation Study Protocol
% LaTeX source - compile with pdflatex
\documentclass[runningheads,a4paper]{llncs}
\usepackage[T1]{fontenc}
\usepackage{booktabs}
\usepackage{array}
\usepackage{multirow}
\usepackage{enumitem}
\usepackage{xcolor}
\usepackage[margin=2.5cm]{geometry}

\definecolor{accent}{HTML}{1a1a2e}
\definecolor{hdrbg}{HTML}{2d2d5e}

\begin{document}

\title{NCD-CIE Clinician Evaluation Study Protocol}
\subtitle{Evaluating LLM Augmentation Layers for Clinical Decision Support}
\author{NCD-CIE Research Team}
\institute{February 2026}
\maketitle

\section{Objective}

Evaluate NL parser accuracy and explanation quality of NCD-CIE's three LLM augmentation layers with practicing clinicians. This study employs a \textbf{mixed-methods design} combining quantitative task accuracy measurements, Likert-scale quality ratings, and qualitative feedback from domain experts.

The study is structured in three parts: (A) automated NL parser accuracy benchmarking, (B) clinician-rated explanation quality assessment, and (C) an optional blinded comparative evaluation against expert-written explanations.

\section{Part A: NL Parser Accuracy}

\emph{No clinicians needed --- fully automated evaluation.}

\begin{itemize}[leftmargin=*]
  \item \textbf{Participants:} 3 LLM backends --- GPT-4, Claude, Ollama (local)
  \item \textbf{Materials:} 50-query test suite (already developed), covering risk queries, what-if scenarios, demographic filters, and multi-NCD composites
  \item \textbf{Protocol:} Run each query through the NL parser with each LLM backend. Score against ground truth: 1.0 (full match), 0.5 (partial match), 0.0 (incorrect). Report per-category accuracy and failure mode analysis.
  \item \textbf{Deliverable:} Table of parse accuracy by query type $\times$ LLM backend
\end{itemize}

\section{Part B: Explanation Quality Assessment}

\begin{itemize}[leftmargin=*]
  \item \textbf{Participants:} 10--15 clinicians (physicians, nurses, pharmacists)
  \item \textbf{Recruitment:} University hospital colleagues, medical faculty
  \item \textbf{Materials:} 20 patient scenarios (diverse age, sex, risk levels, multi-NCD). Each scenario generates: risk assessment, what-if simulation, LLM explanation, and counterfactual narrative.
\end{itemize}

\noindent\textbf{Protocol per scenario (5--7 min):} Read patient profile $\rightarrow$ review system output + explanation $\rightarrow$ rate on a 5-point Likert scale across 6 dimensions:

\begin{table}[h]
\centering
\caption{Likert Rating Dimensions}
\begin{tabular}{ll}
\toprule
\textbf{Dimension} & \textbf{Statement} \\
\midrule
Accuracy & ``The risk explanation is medically accurate'' \\
Clarity & ``The explanation is easy to understand'' \\
Usefulness & ``This would help me in clinical decision-making'' \\
Trustworthiness & ``I would trust sharing this with a patient'' \\
Completeness & ``The explanation covers the key factors'' \\
Actionability & ``The what-if result provides actionable guidance'' \\
\bottomrule
\end{tabular}
\end{table}

\begin{itemize}[leftmargin=*]
  \item \textbf{Open-ended:} ``What would you change about this explanation?''
  \item \textbf{Time per clinician:} $\sim$2 hours (20 scenarios $\times$ $\sim$6 min each)
\end{itemize}

\section{Part C: Comparative Evaluation (Optional)}

\begin{itemize}[leftmargin=*]
  \item \textbf{Design:} Blinded comparison for 10 selected scenarios
  \item \textbf{Procedure:} Present TWO explanations in randomised order: (A) NCD-CIE LLM-generated, (B) Expert-written (cardiologist)
  \item \textbf{Task:} Clinician picks preferred explanation and rates both on the same Likert dimensions
  \item \textbf{Deliverable:} Win rate, preference statistics, paired $t$-test
\end{itemize}

\section{Statistical Analysis}

\begin{table}[h]
\centering
\caption{Statistical Analysis Plan}
\begin{tabular}{lll}
\toprule
\textbf{Analysis} & \textbf{Method} & \textbf{Criterion / Note} \\
\midrule
Inter-rater reliability & Krippendorff's $\alpha$ (ordinal Likert) & $\alpha \geq 0.67$ acceptable \\
Overall quality & Mean $\pm$ SD per dimension, 95\% CI & Target: mean $\geq$ 4.0/5.0 \\
Comparison A vs B & Wilcoxon signed-rank test (paired) & $p < 0.05$ \\
Subgroup effects & By clinician role (MD vs nurse) & Kruskal--Wallis \\
Sample size justification & $N=12$ detects $d=0.8$ at $\alpha=0.05$ & Power $= 0.80$ \\
\bottomrule
\end{tabular}
\end{table}

\section{Patient Scenarios}

Twenty synthetic patient scenarios covering diverse demographics, risk profiles, and intervention types:

\begin{table}[h]
\centering
\caption{Patient Scenario Overview}
\footnotesize
\begin{tabular}{clllll}
\toprule
\textbf{\#} & \textbf{Age} & \textbf{Sex} & \textbf{Key Risk Factors} & \textbf{Scenario Type} & \textbf{Intervention} \\
\midrule
1  & 55 & M & High LDL + smoker          & Single intervention & Statin \\
2  & 62 & F & T2DM + obese (BMI 34)      & Multi-intervention  & Weight + exercise \\
3  & 48 & M & Hypertension (SBP 165)     & What-if             & SBP reduction $-15$ mmHg \\
4  & 70 & F & CKD Stage 3 + CVD history  & Multi-NCD composite & Composite risk \\
5  & 35 & M & Low risk + family history   & Reassurance         & Monitoring \\
6  & 58 & M & Diabetic + hypertensive    & Dual medication     & SGLT2i + BP med \\
7  & 45 & F & Overweight + sedentary     & Lifestyle only      & Exercise + diet \\
8  & 67 & M & Post-MI + high LDL         & Targeted therapy    & PCSK9i \\
9  & 52 & F & Pre-diabetic (HbA1c 6.2\%) & Preventive          & Lifestyle \\
10 & 73 & M & CVD + T2DM + CKD           & Complex multi-drug  & Multi-drug \\
11 & 41 & F & Smoker + oral contraceptives & Cessation         & Smoking cessation \\
12 & 60 & M & Heavy drinker + HTN        & Reduction           & Alcohol reduction \\
13 & 55 & F & Obese + sleep apnea        & Surgical            & Weight loss surgery \\
14 & 49 & M & Athlete + family hx CAD    & Exercise paradox    & Risk stratification \\
15 & 65 & F & Well-controlled T2DM       & De-escalation       & Medication review \\
16 & 78 & M & Frail + polypharmacy       & Age-specific        & Deprescribing \\
17 & 38 & F & Gestational DM history     & Long-term risk      & T2DM prevention \\
18 & 56 & M & Metabolic syndrome (5/5)   & Multi-pathway       & Combined intervention \\
19 & 44 & F & Isolated high TG           & Targeted lipid      & Fibrate/omega-3 \\
20 & 71 & M & eGFR declining (55$\to$48) & Progression prev.   & CKD management \\
\bottomrule
\end{tabular}
\end{table}

\section{Ethics \& IRB}

\begin{itemize}[leftmargin=*]
  \item \textbf{IRB approval required} --- submit to university ethics committee
  \item Written informed consent from all participating clinicians
  \item No real patient data used (all synthetic scenarios)
  \item Classification: Likely ``exempt'' or ``expedited'' (no patient involvement)
  \item Expected IRB timeline: 2--4 weeks
\end{itemize}

\section{Timeline}

\begin{table}[h]
\centering
\caption{Study Timeline}
\begin{tabular}{ll}
\toprule
\textbf{Week} & \textbf{Activity} \\
\midrule
Week 1   & Prepare materials, submit IRB application \\
Week 2--3 & IRB review period \\
Week 3   & Run Part A: NL parser evaluation (no IRB needed) \\
Week 4   & Recruit clinicians, schedule evaluation sessions \\
Week 5   & Data collection (3--4 clinicians/day) \\
Week 6   & Statistical analysis + write-up \\
\bottomrule
\end{tabular}
\end{table}

\section{Expected Impact on Paper}

\begin{itemize}[leftmargin=*]
  \item Adds \textbf{Table X} to Section 7 showing clinician quality ratings across all 6 dimensions
  \item Expected to push peer reviewer score from \textbf{6.5 $\to$ 7.5+}
  \item Closes the last major gap identified across 4 review rounds
  \item Significantly strengthens the ``LLM-Augmented'' contribution claim
  \item Provides empirical evidence of clinical utility beyond automated metrics
\end{itemize}

\end{document}
